% \clearpage
% \vskip -0.2cm
\section{\ours}

\textbf{Framework Overview.}
As shown in Fig.\ref{fig:method}, {\ours} begins by extracting structured parts from the raw-text idea, followed by a heterogeneous deep-knowledge search engine that iteratively acquires and filters high-quality background knowledge from diverse online sources.
The retrieved results, after fine-grained grounding, serve as key references for multi-dimensional and multi-perspective evaluation, empowered by a group of dimension-specialized evaluator agents and an innovation review board composed of diverse academic personas.
Finally, an actionable evaluation report will be synthesized, including comprehensive meta-reviews and concrete suggestions for future revision.
All the notations in this section are summarized in Tab.\ref{tab:notation}.

% \vskip -0.2cm
\subsection{Heterogeneous Deep Knowledge Search}
\label{sec:search_engine}
To maximize usability, we impose no constraints on the format of an idea to be evaluated.
However, to enable precise and efficient knowledge searching and grounding, we first distill the raw textual idea $\boldsymbol{T}$ into the structured representation defined in \S\ref{sec:definition}.
This process is carried out by an \textbf{Extraction Agent} $\boldsymbol{\mathcal{M}}_e$, formally defined as $\boldsymbol{\mathcal{I}}=\boldsymbol{\mathcal{M}}_e(\boldsymbol{T})$.
Then, an exhaustive \textbf{Search Agent} $\boldsymbol{\mathcal{M}}_s$ is deployed to systematically retrieve academic resources relevant to the idea.
To ensure \textit{comprehensiveness}, we expose our research agent to heterogeneous online sources spanning scholarly literature (\arxiv arXiv, \semantic Semantic Scholar, \scholar Google Scholar), web content (\google Google Search), and open-source code repositories (\github Github, \kaggle Kaggle).
By relying solely on online sources, we avoid the staleness inherent to static local corpora and retain \textit{flexibility} to incorporate emerging knowledge.
Moreover, timestamps can be effortlessly integrated into online searching APIs, allowing us to explicitly account for the \textit{timeliness}, a critical attribute of all research ideas.
To ensure retrieval \textit{quality}, we employ an iterative search strategy that simultaneously scales breadth and drills for depth.
We further integrate a fast-slow search hybrid into the deep iterative process to improve search \textit{efficiency}.

\textbf{Fast Knowledge Search.}
Concretely, for every component $\boldsymbol{p}\in \boldsymbol{\mathcal{I}}$ and every search tool $\boldsymbol{u}\in$ (\arxiv,\semantic,\scholar,\google,\github,\kaggle), the search agent first generates a group of tailored queries and then enriches them with synonym expansions, an essential step given that equivalent concepts are routinely phrased differently across communities, yielding the final query set $\boldsymbol{\mathcal{Q}}_{p,u}=\boldsymbol{\mathcal{M}}_s(\boldsymbol{p},\boldsymbol{u})$.
We then directly use the corresponding API for fast retrieval to obtain brief search results:
% \vskip -0.1cm
\begin{align}
\label{eqn:search}
    \boldsymbol{\widetilde{\mathcal{K}}}_{p,u}=\boldsymbol{u}(\boldsymbol{\mathcal{Q}}_{p,u},\boldsymbol{t}).
\end{align}
We partition all the briefly searched knowledge $\boldsymbol{\widetilde{\mathcal{K}}}$ into three types.
Literature knowledge $\boldsymbol{\widetilde{\mathcal{K}}}_l$ contains papers' meta-info such as title, abstract, venue, etc.
Web $\boldsymbol{\widetilde{\mathcal{K}}}_w$ and code knowledge $\boldsymbol{\widetilde{\mathcal{K}}}_c$ include the snapshots of the webpages or repos.
We introduce a \textbf{timestamp} $\boldsymbol{t}$ during search to separately retrieve and process knowledge published before and after $\boldsymbol{t}$.
Pre-knowledge is used to evaluate the idea, whereas post-knowldege are used to suggest revisions.
\textbf{But for notational simplicity, we do not distinguish them symbolically}.

\textbf{Knowledge Ranking and Filtering.}
After discarding unusable knowledge with missing information or inaccessible links, we rank all remaining knowledge using a hybrid scoring function.
Then we calculate the semantic similarity between the idea $\boldsymbol{\mathcal{I}}$ and each piece of knowledge with an embedding model and retain the top $3m$ items for every knowledge type.
A reranking model refines these $3m$ candidates and produces the reranked scores $\boldsymbol{\mathcal{S}}^{\textrm{sem}}$.
We condense the above procedure into the following formulation:
% \vskip -0.1cm
\begin{align}
\label{eqn:semantic_score}
    \boldsymbol{\mathcal{S}}^{\textrm{sem}}_j=\textbf{\texttt{Rerank}}(\textbf{\texttt{Embed}}(\boldsymbol{\mathcal{I}},\boldsymbol{\widetilde{\mathcal{K}}}_j)), j\in\{l,w,c\}.
\end{align}
We then use $\boldsymbol{\mathcal{M}}_s$ as a judge model to assess the relevance and quality of each knowledge, incorporating factors such as paper citations, published venue, website popularity, repo stars, etc., and obtain the model-as-judge scores:
% \vskip -0.1cm
\begin{align}
\label{eqn:llm_score}
    \boldsymbol{\mathcal{S}}^{\textrm{llm}}_j=\boldsymbol{\mathcal{M}}_s(\boldsymbol{\mathcal{I}},\boldsymbol{\widetilde{\mathcal{K}}}_j),\quad j\in\{l,w,c\}.
\end{align}
Finally, we re-weight the two scores with a coefficient $\alpha$ and select the top-$m$ knowledge pieces for each category $j$:
% \vskip -0.1cm
\begin{align}
\label{eqn:hybrid_score}
    \boldsymbol{\widetilde{\mathcal{K}}}^*_j=\textbf{\texttt{Top}}_m(\boldsymbol{\widetilde{\mathcal{K}}}_j, \boldsymbol{\mathcal{S}}_j),\ \boldsymbol{\mathcal{S}}_j=\alpha\boldsymbol{\mathcal{S}}^{\textrm{sem}}_j+(1-\alpha)\boldsymbol{\mathcal{S}}^{\textrm{llm}}_j.
\end{align}
We adopt a hybrid scoring function to simultaneously mitigate the fragility of semantic similarity, and the bias and hallucination inherent in model-as-judge evaluation.

\textbf{Slow Knowledge Search.}
After obtaining the carefully filtered knowledge, we apply a slow search to enrich their content.
For literature-type, we retrieve the full PDF, convert it into raw text, and process it with $\boldsymbol{\mathcal{M}}_e$ to yield structured text that serves as the enriched knowledge $\boldsymbol{\mathcal{K}}_l$.
For web-type, we fetch the webpage via its URL, convert it into raw text, and summarize it into a report with $\boldsymbol{\mathcal{M}}_s$ as $\boldsymbol{\mathcal{K}}_w$.
For code-type, we meticulously search within the repository to collect file-level and function-level call graphs, analyze core code snippets, then combine them with the \texttt{README.md} file to ensemble into a report via $\boldsymbol{\mathcal{M}}_s$ as $\boldsymbol{\mathcal{K}}_c$.

\textbf{Iterative Query Refinement and Search.}
To further uncover latent background knowledge, we task $\boldsymbol{\mathcal{M}}_s$ with refining the original query based on the enriched knowledge, operating along three axes:
\textit{i)} rewriting queries that yielded insufficiently relevant results;
\textit{ii)} generalizing queries that returned overly specific results;
\textit{iii)} concretizing queries that retrieved overly broad results, formulating as:
% \vskip -0.1cm
\begin{align}
\label{eqn:query_refine}
    \widehat{\boldsymbol{\mathcal{Q}}}_{p,u}=\boldsymbol{\mathcal{M}}_s(\boldsymbol{p},\boldsymbol{u},\boldsymbol{\mathcal{Q}}_{p,u},\boldsymbol{\mathcal{K}}_{p,u}).
\end{align}
$\boldsymbol{\mathcal{K}}_{p,u}$ denotes the enriched knowledge retrieved using idea part $\boldsymbol{p}$ and search tool $\boldsymbol{u}$.
We then repeat the search of Eqn.\ref{eqn:search} with the refined query set $\widehat{\boldsymbol{\mathcal{Q}}}_{p,u}$, merge the retrieved brief knowledge $\boldsymbol{\widetilde{\mathcal{K}}}_j$ with the knowledge $\boldsymbol{\widetilde{\mathcal{K}}}^*_j$ retained from the previous round, rank and filter the union $\boldsymbol{\widetilde{\mathcal{K}}}_j\cup\boldsymbol{\widetilde{\mathcal{K}}}^*_j$ via Eqn.\ref{eqn:semantic_score}-\ref{eqn:hybrid_score}, enrich the final top-$m$ pieces, and generate a new batch of queries through Eqn.\ref{eqn:query_refine}. The entire procedure is iterated $\boldsymbol{N}$ times, yielding the final enriched knowledge set $\boldsymbol{\mathcal{K}}$.

% \vskip -0.2cm
\subsection{Knowledge Grounding}
\label{sec:knowledge_grounding}
For a more well-founded evaluation, we further employ a \textbf{Grounding Agent} $\boldsymbol{\mathcal{M}}_g$ to align ideas with knowledge at a finer granularity by excavating the specific connections between the idea $\boldsymbol{\mathcal{I}}$ and the retrieved knowledge.
For each $\boldsymbol{p}\in\boldsymbol{\mathcal{I}}$, we collect all knowledge $\boldsymbol{\mathcal{K}}_p$ retrieved based on it.
Then, for every $\boldsymbol{k}_p\in\boldsymbol{\mathcal{K}}_p$, we distill the most relevant evidences $\boldsymbol{e}_p$ that genuinely supports or contradict $\boldsymbol{p}$ from $\boldsymbol{k}_p$ and provide a detailed relevance analysis $\boldsymbol{s}_p$:
% \vskip -0.1cm
\begin{align}
    \boldsymbol{e}_p,\boldsymbol{s}_p=\boldsymbol{\mathcal{M}}_g(\boldsymbol{p},\boldsymbol{k}_p),\quad\boldsymbol{k}_p\in\boldsymbol{\mathcal{K}}_p.
\end{align}
Finally, we present the grounding results $\boldsymbol{\mathcal{G}}$ fed into the subsequent evaluation module as:
\begin{align}
    \boldsymbol{\mathcal{G}}=\{(\boldsymbol{p},\boldsymbol{\mathcal{G}}_p)\}_{p\in\boldsymbol{\mathcal{I}}},\quad \boldsymbol{\mathcal{G}}_p=\{(\boldsymbol{e}_p,\boldsymbol{s}_p)^i\}^{|\boldsymbol{\mathcal{K}}_p|}_{i=1}
\end{align}

% \vskip -0.2cm
\subsection{Multi-dimensional Multi-perspective Evaluation}
\label{sec:evaluation}
\textbf{Innovation Review Board.}
To achieve review consensus, we augment our evaluation module with a carefully curated innovation review board $\boldsymbol{\mathcal{P}}$ specially designed for academic review.
Each $\boldsymbol{\rho}\in\boldsymbol{\mathcal{P}}$ is an academic persona that comprises an academic profile, a knowledge familiarity vector indicating the degree of acquaintance with literature, web, and code knowledge, and reviewing habits.
During evaluation, we randomly mask a proportion of knowledge corresponding to the persona's familiarity levels across different sources, thereby reflecting actual command of background knowledge.
Other details about the construction, utilization, and examples of the review board are provided in Appx.\ref{app:personas}.

\textbf{Multi-dimensional Decoupled Evaluation.}
We provide five initial evaluation dimensions: \textit{Clarity}, \textit{Validity}, \textit{Novelty}, \textit{Feasibility}, and \textit{Significance}.
Users can also easily register any new evaluation dimensions according to their own needs.
Given the evaluation criteria set $\boldsymbol{\Psi}$, for each $\boldsymbol{\psi}\in\boldsymbol{\Psi}$, we employ a dedicated agent evaluator $\boldsymbol{\mathcal{M}}_{\psi}$ to assess the idea.
To more faithfully emulate human peer review and improve evaluation efficiency, instead of employing all personas in $\boldsymbol{\mathcal{P}}$, we randomly assign five distinct personas as reviewers to each idea.
Given the selected subset $\boldsymbol{\mathcal{P}}'$, for each $\boldsymbol{\rho}\in\boldsymbol{\mathcal{P}}'$ and $\boldsymbol{\psi}\in\boldsymbol{\Psi}$, we score the idea in $[0,10]$ with a detailed review narrative that explains the underlying reasoning:
% \vskip -0.1cm
\begin{align}
    \boldsymbol{\varphi}_{\rho,\psi}=\boldsymbol{\mathcal{M}}_{\psi}(\boldsymbol{\rho},\boldsymbol{\mathcal{I}},\boldsymbol{\mathcal{G}}).
\end{align}
We devise detailed evaluation guidelines and scoring rubrics for each reviewing agent to ensure rigorous.

% \vskip -0.2cm
\subsection{Report Generation}
\label{sec:report}
\textbf{Point-wise Report Synthesis.}
After evaluation, we synthesize all historical information to produce a comprehensive review report for the idea as formalized in Eqn.\ref{eqn:point_report}.
We directly adopt the enriched knowledge reports $\boldsymbol{\mathcal{K}}$ as the background knowledge of the idea.
For revision suggestions $\boldsymbol{\mathcal{V}}$, the \textbf{Report Agent} $\boldsymbol{\mathcal{M}}_r$ proposes future improvements by combining the retrieved future information about the idea with its internal knowledge: $\boldsymbol{\mathcal{V}}=\boldsymbol{\mathcal{M}}_r(\boldsymbol{\mathcal{I}},\boldsymbol{\mathcal{G}}_\textrm{future})$.
Further, we integrate all evaluation information $\{\boldsymbol{\varphi}_{\rho,\psi}\}_{\rho\in\boldsymbol{\mathcal{P}}',\psi\in\boldsymbol{\Psi}}$ and instruct the report agent to produce a meta-review $\boldsymbol{\varphi}_\textrm{meta}$ that contains a detailed summary of the evaluation, a final score $\boldsymbol{s}_\textrm{point}$, and a final decision $\boldsymbol{d}_\textrm{point}$.
Thus, the evaluation conclusion $\boldsymbol{E}_\textrm{point}$ for this idea is expressed as follows:
% \vskip -0.1cm
\begin{align}
    \boldsymbol{E}_\textrm{point}=\{\{\boldsymbol{\varphi}_{\rho,\psi}\}_{\rho\in\boldsymbol{\mathcal{P}}',\psi\in\boldsymbol{\Psi}},\boldsymbol{\varphi}_\textrm{meta}\}.
\end{align}
The final point-wise idea report is $\boldsymbol{P}_{\textrm{point}}=\{\boldsymbol{\mathcal{K}},\boldsymbol{\mathcal{V}},\boldsymbol{E}_{\textrm{point}}\}$.

\textbf{Group-wise Report Synthesis.}
Given a group of $n(n\ge2)$ ideas $\{\boldsymbol{\mathcal{I}}_i\}_{i=1}^n$, we first synthesize point-wise reports $\boldsymbol{P}^{\boldsymbol{\mathcal{I}}_i}_\textrm{point}$ for each idea.
We then feed the reports of all ideas $\{\boldsymbol{P}^{\mathcal{I}_i}_{\textrm{point}}\}_{i=1}^n$ directly as input and enable the report agent to compare every idea along the five aspects in $\boldsymbol{\Psi}$, providing a detailed reasoning process and ultimately producing a comprehensive ranking of all ideas:
% \vskip -0.1cm
\begin{align}
    \boldsymbol{E}_\textrm{group}=\{\{\boldsymbol{\varphi}^\textrm{group}_{\psi}\}_{\psi\in\boldsymbol{\Psi}},\boldsymbol{\varphi}^\textrm{group}_\textrm{meta}\},
\end{align}
where $\boldsymbol{\varphi}^\textrm{group}_{\psi}$ denotes the comparative analysis of different ideas under the specific dimension $\boldsymbol{\psi}$, and $\boldsymbol{\varphi}^\textrm{group}_\textrm{meta}$ contains the final comparative analysis together with the ranked list of ideas.
The final report is $\boldsymbol{P}_{\textrm{group}}=\{\{\boldsymbol{P}^{\mathcal{I}_i}_{\textrm{point}}\}_{i=1}^n, \boldsymbol{E}_{\textrm{group}}\}$.
